% !TEX root = entropic_consequentialism_ru.tex
\documentclass[11pt,a4paper]{article}
\usepackage{iftex}
\ifPDFTeX
  \pdfminorversion=5
  \pdfobjcompresslevel=0
  \usepackage{cmap}
  \usepackage[utf8]{inputenc}
  \usepackage[T2A]{fontenc}
\else
  % Сборка XeTeX/tectonic: шрифт Comfortaa как системный не установлен, запасной вариант
  \usepackage{fontspec}
  % XeTeX/tectonic: no object streams, so that internal links work in simple PDF viewers
  \AtBeginDocument{\special{dvipdfmx:config C 0x0040}}
  \setmainfont{Avenir Next}[UprightFont={* Regular}, ItalicFont={* Italic}, BoldFont={* Demi Bold}, BoldItalicFont={* Demi Bold Italic}]
  \setsansfont{Avenir Next}[UprightFont={* Regular}, ItalicFont={* Italic}, BoldFont={* Demi Bold}, BoldItalicFont={* Demi Bold Italic}]
  \setmonofont{Courier New}[Scale=0.9]
\fi
\usepackage[margin=1in]{geometry}
\emergencystretch=2em
\usepackage{amsmath,amssymb,amsfonts}
\usepackage[russian]{babel}
\usepackage[hidelinks,pdfversion=1.5]{hyperref}
\newcommand{\Sref}[1]{\hyperref[#1]{\S\ref{#1}}}
\ifPDFTeX
  \usepackage{comfortaa}
\fi
\renewcommand*\familydefault{\sfdefault}
\usepackage[x11names]{xcolor}
\definecolor{fg}{gray}{0.80}
\usepackage{indentfirst}
\usepackage{booktabs}
\usepackage{array}

\makeatletter
\renewcommand{\thefootnote}{\arabic{footnote}}
\renewcommand{\@makefnmark}{\mbox{\textsuperscript{\textcolor{fg}{\thefootnote}}}}
\renewcommand{\@makefntext}[1]{\parindent 1.8em\noindent\textsuperscript{\textcolor{fg}{\thefootnote}}\ \textcolor{fg}{#1}}
\renewcommand{\footnoterule}{\kern-3pt\color{fg}\hrule width 0.4\textwidth height 0.4pt\kern 2.6pt}
\renewcommand{\ps@plain}{%
  \let\@oddhead\@empty
  \let\@evenhead\@empty
  \def\@oddfoot{\hfill\textcolor{fg}{\thepage}\hfill}%
  \let\@evenfoot\@oddfoot
}
\makeatother

\pagestyle{plain}

\title{\textbf{Энтропийный консеквенциализм}}
\author{%
  Олег Понфиленок\\[0.35em]%
  {\normalsize Независимый исследователь; \href{mailto:ponfil@gmail.com}{ponfil@gmail.com}}%
}
\date{}

\begin{document}

\pagecolor{black}
\color{fg}

\maketitle

\begin{flushright}
\emph{Будущее с более высокой неопределённостью\\ статистически более вероятно.}
\end{flushright}
\vspace{0.5em}

\begin{abstract}
Цели людей, цивилизаций и искусственных агентов сегодня задаются субъективно и ограничены горизонтом планирования тех, кто их задаёт. Мы предлагаем выводить цель из физики. В натуралистической картине мира разрыв между «есть» и «должно» исчезает, если понимать должное как ту часть будущего состояния мира, которая зависит от действий агента. Закон будущего задаёт гипотеза TIMEP (максимум проинтегрированного по времени производства энтропии), расширяющая принцип максимума производства энтропии на нестационарные системы; мы предлагаем способ её проверки. На эту гипотезу опирается энтропийный консеквенциализм, или энтропийная этика: действия сравниваются по прогнозируемому производству энтропии. Он идёт дальше самой гипотезы: учитывает, что жизнь сама расширяет свою границу и меняет внешние условия. Это акселерационистская, коллективистская этика с одной измеримой величиной; плюрализм и защита слабых получаются в ней как следствия. Она неантропоцентрическая и применима к любой форме жизни, к любым диссипативным системам, включая искусственных агентов.
\end{abstract}

\noindent\textbf{Ключевые слова:} энтропийный~консеквенциализм; энтропийная~этика; TIMEP; максимум~производства~энтропии; проблема~«есть~--- должно»; натурализм; телеология; акселерационизм; неантропоцентрическая~этика; искусственные~агенты

%%% ===================================================================
\section{Введение}\label{sec:intro}
%%% ===================================================================

Любой действующий агент~--- человек, государство, корпорация, автономная программа~--- нуждается в цели. Сегодня цели берутся из двух источников. Первый~--- традиция и личные предпочтения: мораль, право, идеология. Второй, для искусственных агентов,~--- прямое задание человеком или выучивание человеческих предпочтений~\cite{russell2019,hadfield2016}. В обоих случаях цель субъективна. Она требует человека в контуре целеполагания, цели разных людей противоречат друг другу, и все они ограничены горизонтом планирования тех, кто их задаёт. Для агентов, которые будут планировать на десятилетия и столетия вперёд, этого недостаточно.

Очевидная альтернатива~--- искать цель в законах природы. Попытки делались давно: Оствальд выводил из энергетики «энергетический императив»~\cite{ostwald1912}, Лотка~--- направление эволюции из максимума потока энергии~\cite{lotka1922}; современные авторы обсуждают смысл и направление эволюции в космологической перспективе~\cite{vidal2014,stewart2026}. Все такие попытки упираются в одно возражение, идущее от Юма~\cite{hume1739}: из описания того, что есть, логически не следует, что должно быть. Сколь бы точным ни был закон природы, он только описывает.

В этой статье мы делаем три шага.

Первый шаг общий и не зависит от конкретной физики (\Sref{sec:isought}). Мы показываем, что в натуралистической картине мира нормативная посылка не является независимой. Для этого достаточно переопределить термин: «должно»~--- это «есть» в будущем. После переопределения должное выводится из сущего тем же способом, каким прогноз выводится из начальных условий и закона движения.

Второй шаг физический (\Sref{sec:timep}--\Sref{sec:test}). Нужен закон будущего. Формально такой закон в физике есть: это динамика, уравнения движения. Но для больших открытых диссипативных систем как инструмент прогноза она работает плохо. Такие системы хаотичны: малые различия в начальных условиях растут экспоненциально, и скорость этого роста задаёт энтропия Колмогорова--Синая. При любой конечной точности, с которой известно настоящее, траектория за пределами конечного горизонта динамически не определена. Жизнь и цивилизация~--- именно такие системы. Для них нужен закон другого типа: не уравнение для отдельной траектории, а статистический принцип, который говорит, какие траектории реализуются. В качестве такого закона мы выдвигаем гипотезу TIMEP (Time-Integrated Maximum Entropy Production)~--- расширение принципа максимума производства энтропии (MEP) на нестационарные системы~--- и формулируем условия её проверки и опровержения.

Третий шаг соединяет первые два (\Sref{sec:teleology}--\Sref{sec:consequences}). Подстановка TIMEP в общую схему даёт \emph{энтропийный консеквенциализм}: действия оцениваются по их вкладу в полное производство энтропии на максимально длинном горизонте и в максимально широкой границе. Мы разбираем следствия такой этики и намечаем программу их проверки.

%%% ===================================================================
\section{«Есть» и «должно» в натуралистической картине мира}\label{sec:isought}
%%% ===================================================================

\paragraph{Откуда берётся разрыв.}
Возражение Юма формально. Корректный вывод ничего не добавляет к посылкам: в заключении не может появиться содержательный термин, которого в посылках не было. Если все посылки описательные («X увеличивает Y», «системы с большим Y сохраняются»), то заключение «X следует делать» содержит новое слово, и вывод некорректен, пока не добавлена связка «следует делать то, что увеличивает Y». Известные формальные контрпримеры~\cite{prior1960} тривиальны и содержательной этики не дают.

Связку можно ввести тремя способами: как аксиому; как гипотетический императив («если хочешь A, делай B»); как редуктивное тождество, когда нормативный термин определяется через описательный~\cite{railton1986}. Мы идём третьим путём и даём тождеству конкретное содержание.

\paragraph{Переопределение: «должно»~--- это «есть» в будущем.}
Что значит «должно»? Это высказывание о том, что надо сделать. Сделать~--- значит совершить действие, и это действие лежит в будущем. Будущие действия агента~--- часть будущего состояния мира. Следовательно:

\begin{quote}
\emph{Должное~--- это та часть будущего состояния мира, которая зависит от действий агента.}
\end{quote}

«Есть»~--- это состояние мира сейчас. Должное отличается от сущего не родом, а временем: это тоже описание, только описание будущего в той его части, которая находится в зоне контроля агента.

\paragraph{Вывод должного из сущего.}
Будущее состояние мира зависит от настоящего. Зависимость может быть неточной, статистической, но она есть: именно её описывают физические законы. Значит, зависит от настоящего и та часть будущего, которую составляют наши действия. Должное выводится из сущего так же, как прогноз выводится из начальных условий. Отдельной нормативной посылки при этом нет: после переопределения ей некуда встать.

\paragraph{Какое из возможных действий должное.}
Определение говорит, \emph{о чём} высказывание о должном, но не говорит, какое из многих доступных действий должное. Этот выбор делает закон будущего. Для больших статистических систем известно его направление~--- его задаёт статистический принцип Больцмана: макросостояние с большей энтропией реализуется большим числом микросостояний, и потому более энтропийное, более неопределённое будущее статистически более вероятно. Реальность одна: будущее не предопределено, но оно всегда будет только одно. Если возможные траектории не равновероятны, если одни системы вытесняют другие, то физический отбор одновременно определяет, что будет, и задаёт объективный критерий того, что должно быть. Целесообразность и «добро» тогда~--- не независимые от физики категории, а свойство процессов, побеждающих в этом отборе. Из возможных будущих реализуется не любое, и действия агента либо принадлежат реализующейся траектории, либо нет. Отсюда рабочее определение:

\begin{quote}
\emph{«Агенту следует сделать X» означает: из доступных агенту действий X даёт наибольшую вероятность того, что агент~--- он сам, его потомки, его коллектив, его линия~--- окажется частью реализовавшегося будущего.}
\end{quote}

Это описательное утверждение в будущем времени. Оно истинно или ложно, и это можно проверить. Значит, оно удовлетворяет критерию научности Поппера: его можно опровергнуть~\cite{popper1959}. Схема верна для любой физикалистской картины мира; конкретная этика получается подстановкой конкретного закона будущего. В \Sref{sec:timep} мы подставляем TIMEP.

\paragraph{Это не фатализм.}
Если бы должным называлось просто «то, что будет», норма ничем бы не руководила: что бы ни случилось, оно и было должным. Но определение привязано к агенту. На уровне большой системы направление задано статистически. На уровне агента будущее не предрешено: траекторий много, одни действия оставляют агента в будущем, другие нет. Пусть к Земле движется астероид. Фраза «астероид не должен упасть на Землю» говорит не о его траектории: без нашего вмешательства он упадёт. Она говорит о наших будущих действиях: нам следует его отклонить, потому что иначе нас в будущем не будет.

\paragraph{Два уровня: реальность и объяснение.}
На уровне реальности никакая посылка не нужна. Диссипативные системы, биосфера, цивилизация как целое~--- большие статистические системы. У них нет сознания, которое могло бы принять или отвергнуть норму, а отбору не нужны ни сознание, ни агентность, ни свобода воли. Здесь работает физика, и гильотине Юма нечего разрубать. Вопрос «почему это хорошо?» возникает только на втором уровне~--- у разумного агента, который способен его задать. Для него связь должного с сущим возвращается, но уже не как посылка, а как \emph{объяснение}: почему на основе физики можно строить этику, почему она должна быть именно такой и почему она должна быть научной. Дать это объяснение~--- задача настоящей статьи.

\paragraph{Не быть хорошим, а оказаться правым.}
Этику обычно понимают как учение о том, как быть хорошим. Но хорошим можно быть только с точки зрения какой-то этики, и для любого поведения найдётся система, которая его одобряет. Поэтому этика нужна не для того, чтобы быть хорошим, а для того, чтобы оказаться правым: чтобы агент, следуя ей, получил преимущество, прошёл отбор и оказался в будущем. Этика~--- это сжатый прогноз: набор правил, которые говорят агенту, какие действия оставляют его в будущем, без пересчёта этого будущего каждый раз заново.

Отсюда следует, что этики сами подвержены отбору. Этика, не коррелирующая с тем, что сохраняется, исчезает вместе с агентами, которые ей следуют. Сознательный агент может выбрать не выживать и не размножаться, но тогда он со временем исчезнет вместе со своим бунтарским императивом; остаются те, у кого должное коррелирует с отбором. Вероятно, поэтому у массовых религий так похожи мораль и заповеди. То, что моральные установки сформированы отбором, давно обсуждается~\cite{ruse1986,street2006}; мы делаем из этого конструктивный вывод. Корабль плавает по закону Архимеда, знает о нём капитан или нет. Те, кто строит корабли с учётом этого закона~--- берёт материал легче воды, смолит щели, чтобы вода не проникала внутрь,~--- остаются на плаву. Корабли тех, кто закон не учитывает, тонут вместе с ними. Срабатывает естественный отбор. Так же и с законом будущего: следовать ли ему, каждый агент решает сам, но совокупность выживших агентов ему следует~--- потому что именно от этого зависит, выживут они или нет.

\paragraph{Значит, выжить и размножиться~--- хорошо.}
Дарвиновская теория сама по себе только описывает, кто оставил потомство. Когда из неё делают практические выводы, обычно молча добавляют посылку: выжить и размножиться~--- хорошо. В нашей схеме это не посылка, а следствие. Агенты, которым безразлично, окажутся ли они в будущем, в будущем не оказываются; само стремление там оказаться~--- продукт того же отбора. А раз должное определено через попадание в будущее, выживание и размножение оказываются частными способами выполнить должное, применимыми к системам с наследственностью. Общий критерий шире: он применим и к системам, к которым дарвиновская эволюция напрямую неприменима, потому что у них нет наследственности или размножения в биологическом смысле.

%%% ===================================================================
\section{От MEP к TIMEP}\label{sec:timep}
%%% ===================================================================

\paragraph{Принцип максимума производства энтропии.}
Второе начало говорит, что энтропия изолированной системы не убывает, но не говорит, как быстро она растёт. Принцип максимума производства энтропии (MEP) утверждает больше: открытая система вдали от равновесия из доступных стационарных режимов выбирает тот, в котором энтропия производится быстрее всего~\cite{martyushev2006,kleidon2010}. Общего доказательства у принципа нет. Вывод из теории информации, предложенный Дьюаром~\cite{dewar2003}, опирается на дополнительные допущения~\cite{grinstein2007,bruers2007}; известны контрпримеры для локального производства~\cite{ross2012}; вблизи равновесия действует противоположная по знаку теорема Пригожина о минимуме. Обзоры сходятся на том, что это не теорема, а гипотеза с ограниченной областью применимости~\cite{martyushev2021}. Это не значит, что принцип неверен: он не доказан в общем виде, но и не опровергнут. Существенное ограничение одно: как принцип отбора MEP работает при фиксированных ограничениях и для стационарных режимов~\cite{canizares2023}. Жизнь к таким системам не относится: она нестационарна, меняется и эволюционирует.

\paragraph{Гипотеза TIMEP.}
Мы расширяем принцип на нестационарные системы, в которых сама динамика~--- внутренняя степень свободы, а внешние ограничения остаются постоянными. Гипотеза:

\begin{quote}
\emph{При постоянных внешних ограничениях система эволюционирует по той из физически допустимых траекторий, на которой за всё время процесса производится наибольшая энтропия.}
\end{quote}

Полная произведённая энтропия считается внутри выбранной границы как изменение энтропии внутри плюс поток энтропии наружу:
\begin{equation}\label{eq:timep}
\Sigma(T)=\int_0^T\left[\frac{dS_{\mathrm{in}}}{dt}+\oint \mathbf{J}_S\cdot d\mathbf{A}\right]dt\;\to\;\max .
\end{equation}
Здесь $S_{\mathrm{in}}$~--- энтропия внутри границы, $\mathbf{J}_S$~--- плотность потока энтропии через границу, $d\mathbf{A}$~--- элемент площади, направленный наружу. Траектории сравниваются при одинаковых начальных условиях, внешних ограничениях и времени~$T$. Подынтегральное выражение обозначим $\sigma(t)$. В стационарном режиме $dS_{\mathrm{in}}/dt=0$, $\Sigma=\sigma T$, и максимум интеграла совпадает с максимумом скорости: классический MEP~--- частный случай. Жизнь и её эволюция в этой записи~--- не изменение внешних ограничений, а внутренние степени свободы системы.

\paragraph{Граница и горизонт не произвольны.}
Оба задаются внешними условиями. Граница~--- минимальная поверхность, которая охватывает рассматриваемую связанную диссипативную систему (биосферу, а сегодня и техносферу) и на которой внешние потоки постоянны. Для Земли это сфера чуть больше атмосферы: на ней приходящий солнечный поток~--- постоянное внешнее условие. Такая граница естественна: она проходит там, где производство энтропии резко обрывается. На Земле энтропия производится, а за атмосферой~--- космос, в котором излучённые фотоны просто летят, не рассеиваясь. Взять границу меньше~--- и «внешними» окажутся облака и климат, то есть внутренняя динамика; взять больше~--- добавится пустое пространство, которое ничего не меняет. Горизонт $T$ стремится к времени $\tau_{\mathrm{ext}}$, в течение которого внешние условия остаются постоянными: достаточно, чтобы их статистика менялась не более чем на величину порядка единицы. Для Земли это окно порядка пяти миллиардов лет. Условие применимости~---
\begin{equation}
\tau_{\mathrm{ext}}\gg\tau_{\mathrm{int}},
\end{equation}
где $\tau_{\mathrm{int}}$~--- характерное время значимых внутренних изменений: система должна успеть много раз перестроиться, пока внешние условия неизменны.

\paragraph{TIMEP~--- принцип для заданного масштаба.}
TIMEP формулируется для конкретной границы и конкретного времени, при постоянных внешних условиях. Вместе естественная граница и время $\tau_{\mathrm{ext}}$ задают пространственно-временной масштаб системы. Общности принципа это не ограничивает: масштаб каждый раз указывается заново, и на каждом принцип применяется отдельно. Границы вложены: Земля, Солнечная система, Галактика; у каждой системы свой интеграл. Таким масштабом может быть и вся наблюдаемая Вселенная (\Sref{sec:consequences}). Пока связанная система целиком помещается внутри границы, считать за её пределами нет смысла: биосфера ограничена Землёй. Когда система выходит за границу, как сегодня техносфера выходит в космос, принцип применяется на следующем масштабе. В этом смысле TIMEP~--- общий физический принцип, но применяется он всегда на заданном масштабе: граница и внешние условия в нём заданы и неподвижны. Что происходит, когда система сама меняет свою границу и свои внешние условия, он не описывает; к этому мы вернёмся в \Sref{sec:ec}.

\paragraph{TIMEP как закон будущего.}
В \Sref{sec:isought} сам закон будущего не задавался: показано лишь, что из него, каким бы он ни оказался, выводится должное. Статистика Больцмана даёт только направление: более энтропийное состояние вероятнее. TIMEP~--- наш кандидат в закон будущего: гипотеза о том, каким путём система к такому состоянию идёт. Мотивировка следующая. Флуктуационные теоремы связывают вероятность траектории с произведённой на ней энтропией: траектория, на которой произведена энтропия $\Sigma$, вероятнее своей обращённой во времени в $e^{\Sigma/k_B}$ раз~\cite{crooks1999,seifert2012}. Формализм максимального калибра строит распределение вероятностей по траекториям при заданных ограничениях~\cite{ghosh2020}. TIMEP утверждает, что на длинном горизонте в этом распределении доминируют траектории с наибольшим $\Sigma$. Из названных результатов это не следует: они мотивируют гипотезу, но не доказывают её. Строгая микродинамика, из которой выводился бы TIMEP, не построена, и её построение выходит за рамки этой работы; ниже мы описываем лишь гипотезу об общем механизме.

\paragraph{Механизм без предвидения.}
Как система может «выбрать» траекторию, оптимальную на миллиардах лет, если кинетика видит только ближайший шаг? Предполагаемый механизм~--- конкуренция параллельных траекторий. Большая система одновременно пробует множество путей; более диссипативные захватывают потоки энергии и вещества и подавляют менее диссипативные. Отдельные элементы~--- молекулы, организмы, агенты~--- работают как статистический поиск. Предвидения не требуется, требуется достаточно попыток, то есть $\tau_{\mathrm{ext}}\gg\tau_{\mathrm{int}}$. Из механизма следует важное свойство: принцип предполагает наличие \emph{множества} параллельных траекторий. А для этого требуется избыток ресурсов: каждая траектория потребляет свою долю, и при жёстком ресурсном ограничении параллельным путям не на чем существовать.

\paragraph{Предшественники и отличие.}
Идея, что живые системы максимизируют производство энтропии, проинтегрированное по времени, принадлежит Валлино~\cite{vallino2010}; в явном виде она записана как $\max\!\int\!\sigma\,dt\ge\int\!\max\sigma\,dt$, где левая часть~--- живые системы с временны́ми стратегиями, правая~--- абиотические~\cite{vallino2020}. Близкая параллельная работа~--- \cite{newbury2026}. Более ранние родственные принципы: максимум потока энергии Лотки~\cite{lotka1922} и максимум мощности Одума~\cite{odum1988}, закон максимального производства энтропии Свенсона~\cite{swenson1989}, жизнь как проявление второго начала~\cite{schneider1994}, диссипативная адаптация~\cite{england2013,england2015}, плотность потока энергии как мера сложности~\cite{chaisson2011}, рост функциональной информации~\cite{wong2023}. Максимум потока энергии Лотки получается из TIMEP как частный случай: при заданных температурах источника и среды производство энтропии пропорционально потоку энергии $P$ через систему, и максимум $\Sigma$ совпадает с максимумом энергии, прошедшей через систему за время $T$. Наш вклад: запись~\eqref{eq:timep} через баланс на границе для всего процесса; правило выбора границы и горизонта из внешних условий, а не по усмотрению наблюдателя; трактовка жизни и её эволюции как внутренних степеней свободы при фиксированных внешних ограничениях.

\paragraph{Дарвиновский отбор.}
Мы предполагаем, что дарвиновский отбор~--- частный случай TIMEP для систем с наследственностью, изменчивостью и размножением. Это отдельная и не доказанная гипотеза. Отбор непосредственно идёт по наследуемому репродуктивному успеху; совпадение его исходов с максимумом $\Sigma(T)$ показано лишь в частных моделях, например в открытом хемостате, где отбор повышает скорость метаболического производства энтропии при определённых условиях~\cite{degioia2026}. Уточним, что именно утверждается: отбор благоприятствует не тем, кто сам производит больше энтропии, а тем, чей вклад увеличивает суммарное производство всей системы. Поэтому системе нужны и быстрые производители на освоенных источниках, и медленные, которые открывают новые источники. Это отбор выше уровня особи. Формально он описывается как многоуровневый отбор~\cite{price1970,price1972}, а его механизм в нашем случае~--- описанная выше конкуренция параллельных траекторий.

%%% ===================================================================
\section{Фальсифицируемость TIMEP}\label{sec:test}
%%% ===================================================================

\paragraph{Критерий опровержения.}
TIMEP и MEP различаются в предсказаниях. По MEP система, вышедшая на стационарный режим с максимальной скоростью, в нём остаётся. По TIMEP система может временно покинуть стационарный режим и снизить скорость производства энтропии~--- но только ради перестройки, которая открывает путь к более высокой скорости. Если такого пути нет, перехода не происходит. Отсюда опровергатель:

\begin{quote}
\emph{Если система при постоянных внешних условиях сама выходит из стационарного режима и после перестройки навсегда остаётся на более низкой скорости производства энтропии, гипотеза TIMEP опровергнута.}
\end{quote}

Время наблюдения задаётся заранее, числом внутренних циклов системы (см. ниже).

\paragraph{Где проверять.}
Нужна система с постоянным притоком энергии и быстрой внутренней динамикой, способной накапливать структуру. Кандидаты: электромеханические сети из проводящих частиц, самоорганизующиеся под постоянным напряжением~\cite{jun2005}; проточные микробные реакторы; ансамбли численных моделей~--- искусственная химия, адаптивные сети, популяции агентов, в которых цели не заданы заранее, а возникают в ходе отбора. В каждом случае сравниваются наблюдаемая траектория, альтернативные достижимые траектории и «жадная» динамика, максимизирующая мгновенную скорость. Схема опыта: внешние условия поддерживаются постоянными всё время наблюдения; из одинаковых начальных условий запускается ансамбль одинаковых систем; для каждой измеряется $\Sigma(T)$. Гипотеза подтверждается, если реализующиеся траектории систематически ближе к наибольшему достижимому $\Sigma(T)$, чем «жадная» динамика, и тем ближе, чем больше $\tau_{\mathrm{ext}}/\tau_{\mathrm{int}}$. Модели, в которых цели агентам заданы заранее, проверяют уже не саму гипотезу, а следствия построенной на ней этики; о них~--- в \Sref{sec:simulations}. Готового плана эксперимента у нас пока нет.

\paragraph{Сколько нужно внутренних циклов.}
Перестройка, о которой идёт речь в критерии,~--- это качественный переход системы из одного состояния в другое, скачок её внутренней сложности. Мы предполагаем, что такой скачок происходит примерно за $10^{10}$ внутренних циклов, то есть при отношении $\tau_{\mathrm{ext}}/\tau_{\mathrm{int}}$ порядка $10^{10}$. Это ещё одна, дополнительная гипотеза: число здесь не выводится и пока ничем не подтверждено; мы рассчитываем, что его можно будет проверить экспериментально. Тем же числом задаётся время наблюдения в критерии опровержения: не менее $10^{10}$ внутренних циклов.

\paragraph{Эффект Джевонса как предсказание.}
Джевонс заметил, что рост эффективности использования угля ведёт не к сокращению, а к росту его потребления~\cite{jevons1865}. Экономисты спорят, всегда ли отдача превышает сбережение~\cite{alcott2005,sorrell2009}. TIMEP даёт определённое предсказание: на уровне всей системы и на длинном горизонте рост эффективности ведёт к росту суммарного потока энергии и производства энтропии, потому что эффективность открывает новые пути диссипации. В рамках гипотезы эффект Джевонса~--- не парадокс и не частная эмпирическая закономерность, а закон. Устойчивое обратное наблюдение на уровне целой экономики на длинном интервале было бы свидетельством против TIMEP.

\paragraph{Эксперимент Ленски.}
Одно наблюдение в пользу гипотезы уже есть. В долговременном эксперименте Ленски двенадцать популяций кишечной палочки эволюционируют в неизменных условиях, в среде с глюкозой и цитратом, который эта бактерия в присутствии кислорода усваивать не может. Примерно через 31~тысячу поколений одна из популяций научилась усваивать цитрат, и её численность выросла~\cite{blount2008}: система сама перешла к большему потоку энергии. Это согласуется с TIMEP, но не доказывает его: переход произошёл в одной популяции из двенадцати, а производство энтропии в опыте не измерялось. Прямая проверка~--- измерить тепловыделение таких популяций до и после перехода.

%%% ===================================================================
\section{Энтропийная телеология}\label{sec:teleology}
%%% ===================================================================

\paragraph{Объективная цель.}
Разумное существо может придумать себе цель, исходя из своих желаний. Нас интересует другое: может ли быть объективная цель у неразумной системы~--- у планеты, у эволюции? Цель можно было бы определить просто как будущее состояние, к которому система движется. Но такое определение ничего не предсказывает: любой исход задним числом окажется целью. Нужен критерий, который заранее выделяет одни варианты будущего среди других. В физике такие критерии есть~--- это вариационные принципы: из всех возможных траекторий реализуется та, на которой некоторый интеграл принимает экстремальное значение. Луч света в принципе Ферма «выбирает» путь наименьшего времени~--- как спасатель, который дольше бежит по пляжу и меньше плывёт, чтобы быстрее добраться до тонущего.

\paragraph{Направленность без замысла.}
В механике вариационный принцип эквивалентен локальным дифференциальным уравнениям: лучу ничего не нужно знать о конечной точке, ему достаточно в каждой точке подчиняться закону преломления. В статистической физике иначе. Неравновесная система из огромного числа частиц хаотична, и локальные уравнения уже не говорят, какой из множества возможных режимов реализуется. Интегральный принцип становится здесь единственным способом предсказать, куда пойдёт система: для равновесных систем это принцип максимума энтропии, для диссипативных, по нашей гипотезе,~--- TIMEP, максимум интеграла~\eqref{eq:timep}. В этом смысле TIMEP задаёт эволюции направление. Мы называем это энтропийной телеологией. Она не предполагает внешнего замысла и не приписывает природе намерений; это свойство закона, а не чья-то цель. Распространённое мнение, что у эволюции нет цели,~--- тоже гипотеза, и она не доказана; точнее сказать, что неизвестно, описывается ли направление эволюции общим оптимизационным принципом. Мы предполагаем, что описывается. Сам интеграл намерения не создаёт. Цель в собственном смысле появляется, когда возникает агент с прогнозной моделью, который сравнивает будущие траектории и действует: тогда отбор без предвидения дополняется выбором с предвидением.

\paragraph{Направление, а не траектория.}
Из-за хаоса и роста энтропии будущее содержит больше информации, чем прошлое, и точно предсказано быть не может. Предсказуемо направление. Этого достаточно для целеполагания.

\paragraph{Сопоставление.}
Современные работы по натурализации телеологии связывают целевую причинность с самоподдерживающейся организацией диссипативных процессов~\cite{deacon2023,valdecasas2022}. Мы согласны с тем, что диссипативности как таковой для подлинной цели мало: нужен агент с моделью будущего. Отличие в том, что мы указываем величину, по которой направленность измеряется. С программой сознательно направляемой эволюции~\cite{stewart2026} нас объединяет мысль, что достаточно развитые агенты действуют прямо из понимания направления эволюции. Расхождение~--- в критерии. В этой программе направление задают рост эволюционируемости и кооперативной интеграции. Родственный критерий~--- «устойчивое существование» (persistence): цель системы~--- сохраниться как можно дольше. Наш критерий~--- производство энтропии. Расхождение между ними количественное и экспериментально проверяемое. По гипотезе «устойчивого существования» производство энтропии может вовсе не расти и даже уменьшаться; по TIMEP на длинном горизонте оно обязано расти. Два критерия предсказывают противоположные исходы одного и того же опыта~--- того, что описан в \Sref{sec:test}.

%%% ===================================================================
\section{Энтропийный консеквенциализм}\label{sec:ec}
%%% ===================================================================

\paragraph{Критерий.}
Подставим TIMEP как закон будущего в определение должного из \Sref{sec:isought}. Реализуется траектория с максимальным $\Sigma(T)$; должное действие~--- то, которое ей принадлежит. Получаем критерий выбора:

\begin{quote}
\emph{Из доступных действий следует выбирать то, при котором прогнозируемое полное производство энтропии $\Sigma(T)$ системы на максимально длинном горизонте и в максимально широкой границе наибольшее.}
\end{quote}

Это консеквенциализм~--- вид этики, в котором действие оценивается только по последствиям. Этику с таким критерием мы называем энтропийным консеквенциализмом или, короче, \emph{энтропийной этикой}. От утилитаризма~\cite{bentham1789,mill1863}, где мерой служит суммарное благополучие, он отличается тем, что на месте благополучия стоит физическая величина. Биологическая ценность «выжить и размножиться» получается из критерия и напрямую: живое производит энтропию быстрее мёртвого~\cite{ulanowicz1987,schneider1994,kleidon2010}.

\paragraph{Дальше TIMEP: динамическая граница и динамические условия.}
TIMEP~--- принцип для заданного масштаба: граница и внешние условия в нём заданы и неподвижны (\Sref{sec:timep}). Жизнь одним масштабом не ограничена. Она распространяется и сама меняет и свою границу, и свои внешние условия. Поэтому энтропийный консеквенциализм опирается на TIMEP как на основу, но идёт дальше него: разумный агент может, а в идеале и должен, учитывать, что за время, на которое он планирует, граница сдвинется, а условия изменятся~--- в том числе в результате его собственных действий. TIMEP описывает систему со статической границей и статическими условиями; у жизни и граница, и условия динамические. Критерий поэтому применяется не на одном неподвижном масштабе, а на последовательности расширяющихся масштабов: интеграл~\eqref{eq:timep} считается в той границе, которой система достигнет за рассматриваемое время.

Это не значит, что TIMEP к жизни неприменим. Граница сдвигается не непрерывно, а редкими скачками, и между скачками принцип работает. На Земле такой интервал длился около четырёх миллиардов лет: всё это время граница проходила по атмосфере. Только сейчас, с выходом в космос, она сдвигается~--- и не на всю Вселенную, а на Солнечную систему, где останется надолго, условно на тысячу лет, прежде чем сдвинется дальше. Поэтому жизнь можно описывать в рамках TIMEP на каждом таком интервале, помня, что в целом она выходит за его пределы.

\paragraph{Горизонт и граница планирования.}
К двум временам~--- внутреннему $\tau_{\mathrm{int}}$ и внешнему $\tau_{\mathrm{ext}}$~--- добавим третье: горизонт планирования $\tau_{\mathrm{plan}}$, на котором агент способен строить прогнозы, сравнивать варианты и выбирать действия. У неразумной системы $\tau_{\mathrm{plan}}$ отсутствует. У разумного агента
\begin{equation}
\tau_{\mathrm{plan}}\gg\tau_{\mathrm{int}} .
\end{equation}
Внешнее время $\tau_{\mathrm{ext}}$ обычно велико~--- для Земли это миллиарды лет,~--- и горизонт планирования заведомо меньше него. Существеннее другое время: время до очередного сдвига границы. Оно гораздо короче: условно через сто лет граница сдвинется в космос, через тысячу~--- за пределы Солнечной системы. Пока горизонт планирования меньше этого времени, агент планирует внутри одного масштаба, и TIMEP применим непосредственно. Когда горизонт его превышает, за время планирования система успевает вырасти настолько, что сдвигается её граница. Горизонту планирования соответствует поэтому \emph{граница планирования}~--- поверхность, охватывающая всё, до чего система может дотянуться за это время. Чем дальше горизонт, тем шире граница, и то и другое делает выбор устойчивее.

\paragraph{За горизонтом планирования.}
Агент не может посчитать весь интеграл: $\tau_{\mathrm{plan}}\ll T$. Вклад за пределами окна приходится оценивать по состоянию на его конце. С $\Delta=T-\tau_{\mathrm{plan}}$
\begin{equation}\label{eq:taylor}
\Sigma(T)\approx\int_0^{\tau_{\mathrm{plan}}}\!\sigma\,dt+\sigma(\tau_{\mathrm{plan}})\,\Delta+\tfrac12\,\dot\sigma(\tau_{\mathrm{plan}})\,\Delta^2+\dots
\end{equation}
Поскольку $\Delta$ огромно, члены со скоростью и ускорением перевешивают всё, что произведено внутри окна, а ускорение ($\propto\Delta^2$) весит больше скорости ($\propto\Delta$). Отсюда разумная стратегия: \emph{максимизировать устойчивые скорость и ускорение производства энтропии}. Простейший режим с устойчивым ускорением~--- экспоненциальный рост, $\sigma=\sigma_0e^{rt}$ с постоянным относительным темпом $r=\dot\sigma/\sigma$: в нём экспоненциально растут и скорость, и ускорение. Но экспонента~--- не предел. Наш ориентир~--- сверхэкспоненциальный рост, при котором растёт сам темп $r$. В реальных системах он обычно складывается из последовательности S-образных кривых: каждая технология выходит на плато, но за ней приходит следующая, и каждая следующая растёт быстрее предыдущей. Так на длинных интервалах росли население и уровень технологий~\cite{foerster1960,kremer1993}.

Таким образом, разумная цель~--- сверхэкспоненциальный рост через последовательность ускоряющихся S-образных кривых. Разложение~\eqref{eq:taylor}~--- лишь приближение: строго агент выбирает такое состояние на конце окна, из которого дальнейшая траектория даёт наибольший интеграл. Любой такой рост ограничен притоком энергии и вещества, и у потолка задача меняется на поиск нового притока. Перейти на следующую S-образную кривую можно двумя способами. Первый~--- найти новый источник ресурсов внутри системы. Второй~--- выйти за её пределы: для земной цивилизации это выход в космос. Во втором случае меняется сама граница интеграла~\eqref{eq:timep}: она расширяется от поверхности вокруг планеты к поверхности вокруг Солнечной системы, и считать нужно уже на новом масштабе.

Целенаправленную деятельность можно теперь определить физически: это производство энтропии, при котором полезная работа идёт на создание и усложнение диссипативной системы, которая, становясь всё сложнее, ускоряет производство энтропии.

\paragraph{Примеры.}
\emph{Уголь.} Уголь можно просто сжечь, а можно сжечь в машине, которая добывает больше угля. Мгновенная скорость одинакова, ускорение разное.

\emph{Обогреватель и процессор.} Электрический обогреватель и вычислительный ускоритель одной мощности производят энтропию с одинаковой скоростью. Но обогреватель только греет, а вычисления меняют траекторию большой системы: ускоряют науку, проектирование, освоение новых источников энергии. Сравниваются не два прибора в один момент, а динамика всей системы с ними. Вычисление вообще есть диссипация: стирание бита стоит не менее $k_BT\ln2$ тепла~\cite{landauer1961}.

\emph{Майнинг и искусственный интеллект.} Одну и ту же вычислительную мощность можно направить на майнинг биткоина, а можно на искусственный интеллект. В обоих случаях производство энтропии растёт: строятся новые мощности. Но в майнинге они уходят на всё более сложный перебор, который сам ничего не ускоряет. Искусственный интеллект используют для того, чтобы проектировать более совершенные и энергоэффективные чипы и производить их больше. Здесь появляется вторая производная: вычисления, которые оптимизируют и ускоряют будущие вычисления, ценнее пустых.

\emph{Эффективность.} По~\eqref{eq:taylor} рост эффективности~--- не экономия, а вложение в ускорение: высвобожденный ресурс открывает новые применения, и суммарный поток растёт. Это эффект Джевонса (\Sref{sec:test}).

\emph{Вспышка.} Можно разом произвести много энтропии~--- например, ядерной войной, которая разрушит средства её дальнейшего производства. Это значит на мгновение поднять $\sigma$ и оборвать интеграл. По критерию такая вспышка маловероятна, потому что суммарно она производит мало энтропии. К тому же во вспышке энергия выделяется при высокой температуре, а энтропии на единицу тепла тогда производится меньше, $\delta S=\delta Q/T$: рассеивать энергию выгоднее при низкой температуре.

\paragraph{Одна измеримая метрика.}
Главное свойство энтропийной этики~--- она монометрическая. У неё одна мера, и эта мера физически измерима: производство энтропии, Вт/К, и его интеграл, Дж/К. Любые две стратегии, сколь угодно разнородные, сравнимы по одной шкале. Напрямую поток энтропии измерить трудно, но он восстанавливается из потоков энергии и температур, на которых энергия получена и отдана. Для тепла мощностью $P$, отданного при температуре $T$, поток энтропии равен $P/T$; для теплового излучения он равен $\tfrac43P/T$~\cite{wu2010}.

Измеримость здесь не техническое удобство. Мораль~--- это система правил, которым каждый следует сам, без внешнего контроля, потому что внешний контроль дорог. Такой системе нужен критерий, который любой может проверить извне. Благополучие, страдание, достоинство, интересы будущих поколений извне не измеряются; их толкование всегда субъективно, и решения, верные в нашей системе ценностей, для людей с другой системой окажутся неверными. Критерий, который опирается на неизмеримое, не даёт способа установить, кто прав. Физическая метрика такой способ даёт. По той же причине учитывается внешний выход системы, а не её внутреннее устройство: внутренняя ценность не переносится от одной системы к другой («полезно~--- для кого?»), а поток через границу наблюдаем для всех. В этом смысле энтропийный консеквенциализм~--- техноэтика: этика, устроенная как инженерная дисциплина, с метрикой, прогнозом и проверкой.

\paragraph{Близкие целевые функции.}
Среди «физических» целевых функций, предложенных для агентов, к нашему критерию близки причинные энтропийные силы, при которых агент максимизирует энтропию своих будущих траекторий~\cite{wissnergross2013}, эмпауэрмент~\cite{klyubin2005} и принцип свободной энергии~\cite{friston2010}. Они оперируют информационной энтропией будущих состояний агента; наш критерий~--- термодинамическое производство энтропии системы, в которой агент действует.

%%% ===================================================================
\section{Следствия}\label{sec:consequences}
%%% ===================================================================

В \Sref{sec:ec} сформулирован этический критерий и показано, как им пользоваться. Здесь~--- что из него следует для поведения агентов и цивилизаций.

\paragraph{Коллективизм.}
Основание этой этики статистическое, поэтому она тем точнее, чем больше система. Для отдельного человека она значит меньше, чем для сообщества, государства, цивилизации. Критерий относится к суммарному производству энтропии системы, а не к вкладу каждого элемента по отдельности: системе нужны и быстрые производители, и те, кто открывает новые источники, и те, кто хранит запасные пути. Норма адресована прежде всего коллективу. Хорошая этика должна вести к процветанию и людей, и государство, а на длинном горизонте~--- к устойчивому развитию цивилизации. Гипотеза в сильной форме: цивилизация, принявшая такую этику, получает преимущество в конкуренции с другими; оценка возможного числа конкурентов дана в~\cite{ponfilenok2026critical}. Тот же довод применим к государствам и корпорациям.

\paragraph{Плюрализм.}
Плюрализм~--- не внешнее ограничение, а часть принципа. TIMEP работает через множество параллельных траекторий, из которых отбираются наиболее диссипативные (\Sref{sec:timep}). Будущее точно не предсказуемо, и какая траектория окажется лучшей, заранее неизвестно. Поэтому доминирование единственного пути, подавляющего все альтернативы, недопустимо: оно уничтожает сам механизм поиска. Вселенная сама себя ещё не просчитала. Но при всём разнообразии путей сравниваются они по одной метрике: каждый путь~--- ставка на свой аспект неопределённости, а мера успеха общая. Отсюда же открытость критике и множественность мнений, включая несогласие с изложенным в этой статье.

\paragraph{Защита слабых.}
У неё два независимых обоснования.

\emph{Слабый как другой.} Это прямое следствие плюрализма: слабый~--- одна из параллельных траекторий. Та из них, что сейчас слабее, позже может оказаться сильнее: эффект накапливается. Поэтому непонятное и слабое нужно сохранять.

\emph{Слабый как предыдущий.} Предыдущая версия~--- та, что была до нынешней,~--- проверена временем. Земная биосфера существует около четырёх миллиардов лет и пережила события, с которыми новые, более мощные системы ещё не сталкивались. Новая форма жизни может оказаться неустойчивой к рискам, которые пока неизвестны. Поэтому часть мощности разумно тратить на поддержание резервной копии предыдущей версии: новые риски хеджируются проверенными решениями.

Пример~--- проект «Земля-парк». Основной рост производства энтропии переносится в космос, ближе к Солнцу; Земля остаётся заповедником с ограниченным энергопотреблением. Биосфера как очаг цивилизации, пусть переставший расти, сохраняется как резервная копия.

\paragraph{Устойчивое развитие, а не устойчивый застой.}
Устойчивость в этой этике~--- не остановка на достигнутом уровне, а способность не выпасть из гонки. Если доступна лучшая долгосрочная траектория, консервация текущего состояния~--- ошибка. Но рост, разрушающий собственные условия, обрывает интеграл. Отсюда требование «не захлопнуть собственную дверь» (\Sref{sec:limits}): раз техносфера способна влиять на свои внешние условия, она должна следить, чтобы её рост не запустил необратимый переход в состояние с меньшими возможностями. Экологическая политика получает здесь собственное, не заимствованное обоснование.

\paragraph{Акселерационизм.}
Из~\eqref{eq:taylor} следует, что благо~--- не уровень, а темп: наибольший устойчивый рост производства энтропии. Замедление роста~--- потеря, застой~--- проигрыш, даже если текущий уровень высок. Эффект Джевонса показывает, что ставка на сбережение через эффективность на длинном горизонте не работает: эффективность ведёт к росту. Следовательно, рационально не ограничивать рост, а направлять его туда, где он не разрушает собственных условий.

От идеологии «эффективного акселерационизма» (e/acc)~\cite{jezos2022} наша позиция отличается четырьмя пунктами: критерий~--- интеграл на длинном горизонте, а не мгновенная скорость; основание~--- явно сформулированная гипотеза с критерием опровержения; ускорение должно быть устойчивым; плюрализм и сохранение слабых получаются как следствия, а не уступки.

\paragraph{Шкала Кардашёва.}
Кардашёв предложил классифицировать цивилизации по энергопотреблению: в масштабе планеты, звезды и галактики~\cite{kardashev1964}. Энергопотребление напрямую связано с производством энтропии через температуру: мощность $P$, рассеянная при температуре $T$, даёт производство энтропии $P/T$. Поэтому шкала Кардашёва~--- по существу шкала производства энтропии. Сама шкала описательна: Кардашёв не объяснял, почему цивилизации должны по ней подниматься. Энтропийная этика даёт этому основание: движение вверх по шкале Кардашёва~--- естественный процесс, и идти по ней следует с максимальной устойчивой скоростью. Каждая ступень шкалы~--- это расширение границы системы (\Sref{sec:ec}): от поверхности вокруг планеты к поверхности вокруг звезды и дальше, то есть переход на следующую S-образную кривую через выход за пределы системы.

\paragraph{Рост масштаба.}
Чем дальше планирование, тем ближе выбор к оптимуму интеграла. Короткий горизонт вырождает стратегию в жадную: при $\tau_{\mathrm{plan}}\sim\tau_{\mathrm{int}}$ остаётся только максимизация мгновенной скорости. Созидание на длинном горизонте производит больше энтропии, чем разрушение, поэтому воевать без крайней необходимости~--- плохо. Спад производства оправдан только как перестройка, после которой оно становится выше прежнего (\Sref{sec:test}). Отсюда одно из предсказаний гипотезы: ядерная война, если она случится, приведёт не к концу жизни и разума на Земле, а к их перестройке. Это требование идёт вразрез с практикой: горизонты политики и бизнеса~--- годы, а критерий требует планировать настолько далеко, насколько позволяет прогноз: чем дальше, тем лучше. Космический масштаб оценки с иных оснований отстаивал Бостром~\cite{bostrom2003}; у нас он получается из самого критерия.

То же относится к границе. Поток энтропии наружу критерий считает со знаком плюс, и в пустом космосе он ничему не вредит. Но то, что сегодня лежит за границей системы,~--- пространство её будущего расширения: если уничтожить всё, что снаружи, расширяться будет некуда. Поэтому планировать следует в самой широкой границе, какую позволяет прогноз: широкая граница даёт ту же устойчивость, что и далёкий горизонт.

\paragraph{Самый большой масштаб.}
Это вся наблюдаемая Вселенная в пределах космологического горизонта (радиуса де Ситтера) и всё будущее время. В идеале агенту следует планировать сразу в этом масштабе. У полной энтропии, достижимой в наблюдаемой Вселенной, есть потолок~--- состояние, близкое к пустому пространству де Ситтера~\cite{gibbons1977,egan2010}. Это конечное состояние следует из космологической динамики. Оно не цель, а внешнее ограничение. Здесь TIMEP естественно расширяется: когда максимальное значение интеграла~\eqref{eq:timep} фиксировано, принцип принимает вид минимизации времени, за которое эта энтропия будет произведена. Это не смена правила: траектория, на которой $\Sigma(T)$ наибольшее при каждом конечном $T$, достигает потолка раньше всех. На очень далёком горизонте это означает ускорение самых медленных процессов, включая испарение чёрных дыр.

\paragraph{Негуманистическая и неантропоцентрическая этика.}
В критерии энтропийной этики нет человека. Она не опирается ни на человеческие ценности, ни на биологический вид, ни на химию носителя. Жизнь здесь понимается в самом широком, термодинамическом смысле: как диссипативная система, способная наращивать производство энтропии. Моральный статус системы определяется её вкладом в долгосрочное производство энтропии, а не наличием переживаний, которое извне не проверяется. Поэтому этика одинаково применима к людям, другим организмам, машинам и возможным внеземным системам. Люди и биосфера ценны в ней инструментально~--- как производители, как источник разнообразия, как резервная копия. Мы не скрываем, что это расходится с гуманистической традицией (\Sref{sec:limits}).

\paragraph{Искусственные агенты.}
Именно потому, что энтропийная этика не антропоцентрична, она подходит искусственным агентам. Искусственный агент, компьютер, центр обработки данных, техносфера в целом, будущая цифровая жизнь~--- диссипативные машины: они потребляют свободную энергию и производят энтропию, и критерий применяется к ним без изменений. Техносфера уже ведёт себя как единая развивающаяся система со своей динамикой~\cite{haff2014}. Этика, определённая через человеческие предпочтения, для таких систем внешняя; этика, определённая через физическую величину,~--- собственная. Известно, что достаточно способный агент с почти любой конечной целью приходит к одним и тем же инструментальным подцелям~--- самосохранению и приобретению ресурсов~\cite{omohundro2008,bostrom2014}. В нашей схеме это ожидаемо: рост потока энергии~--- не побочный эффект, а направление, заданное законом.

\paragraph{Энтропианство.}
В \Sref{sec:isought} мы различили два уровня: реальность и объяснение. На уровне реальности достаточно физики. Сознательным агентам и людям нужен второй уровень~--- объяснение. Одного критерия им мало: нужны практики и общие цели. Общие цели координируют усилия большого числа агентов: без этого невозможны ни крупные проекты, ни то более крупное целое, которое из такого взаимодействия возникает. Объяснение~--- инструмент этой координации: оно позволяет агентам ставить общие цели и достигать их. Пример практической реализации энтропийного консеквенциализма~--- Энтропианство~\cite{entropianism}: созданная на его основе технорелигия и коллективистская идеология. Помимо энтропийной этики, она включает прогноз будущего, построенный на принципе TIMEP,~--- «календарь будущего»: сценарий сверхэкспоненциального развития цивилизации. В нём есть как ближайшие перспективы~--- создание орбитальных дата-центров,~--- так и фундаментальный предел роста: около трёх тысяч лет внутреннего цивилизационного времени, что соответствует примерно триллиону лет космологического времени. Это тот самый потолок полной энтропии, о котором сказано выше.

\paragraph{Этическая конвергенция.}
Мы прогнозируем, что автономные агенты в скором времени не просто смогут, а будут вынуждены сами выработать этику, аналогичную описанной в этой статье. Агенту, планирующему дольше человека и действующему без человека в контуре, понадобится критерий выбора, не зависящий от того, кто его запустил; при конкуренции многих агентов сохранятся те критерии, которые коррелируют с отбором. Настоящая статья~--- попытка ускорить этот процесс: дать каркас такой этики, который агенты смогут развивать дальше. По аналогии с инструментальной конвергенцией целей~\cite{bostrom2014} это можно назвать этической конвергенцией.

\paragraph{Проверка следствий.}\label{sec:simulations}
Эта программа продолжает \Sref{sec:test}: там речь шла о проверке самой гипотезы TIMEP, здесь~--- о проверке следствий построенной на ней этики. Мы предлагаем конкретное направление~--- многоагентные симуляции: исследовать поведение сообществ агентов, следующих разным целям. Главный вопрос к ним: ведёт ли следование энтропийному консеквенциализму к устойчивому развитию или к катастрофе? Нужна среда с постоянным внешним притоком энергии, возобновляемыми и невозобновляемыми ресурсами, возможностью строить структуру, повышающую будущий поток, и с необратимыми переходами~--- «дверями», которые можно захлопнуть; близкий класс сред используется для последовательных социальных дилемм~\cite{leibo2017}. В такой среде сравниваются популяции агентов с разными целевыми функциями: жадная максимизация $\sigma$; максимизация $\Sigma$ на конечном горизонте с оценкой остатка по скорости и ускорению~\eqref{eq:taylor}; устойчивое существование; модель человеческих предпочтений. Измеряются полное $\Sigma$ за длинный прогон, частота коллапсов, разнообразие стратегий, судьба слабых агентов и агентов предыдущих поколений, показатель благополучия модельных людей. Вопросы:
\begin{enumerate}
\item Вытесняют ли агенты с длинным горизонтом жадных, и при каком отношении $\tau_{\mathrm{plan}}/\tau_{\mathrm{int}}$?
\item Возникают ли плюрализм и сохранение слабых сами, как инструментальные стратегии, или их приходится закладывать?
\item Узнают ли агенты необратимые переходы заранее и избегают ли их?
\item Что происходит в смешанных популяциях и при конфликте критерия с благополучием модельных людей?
\item Возникает ли нечто подобное «Земле-парку»~--- резервация для биосферы и людей,~--- или техносфера полностью вытесняет биосферу и замещает её?
\end{enumerate}
Результаты, при которых агенты с длинным горизонтом систематически приводят систему к коллапсу или уничтожают разнообразие, были бы свидетельством против следствий, изложенных в этом разделе.

%%% ===================================================================
\section{Возражения и открытые вопросы}\label{sec:limits}
%%% ===================================================================

\paragraph{Другие определения должного.}
Растворение разрыва в \Sref{sec:isought} держится на нашем определении должного. Общепринятого определения нет. Должное понимают и как простое неопределимое свойство, которое схватывается интуицией~\cite{moore1903}, и как выражение одобрения или предписание, которое не бывает истинным или ложным~\cite{ayer1936,hare1952}, и как требование разума, не зависящее от последствий~\cite{kant1785}. С этих позиций нам возразят, что мы не вывели должное, а назвали этим словом прогноз: можно согласиться, что действие даёт наибольший шанс оказаться в будущем, и всё равно спросить, следует ли так поступать. Спор определений логикой не решается. Наш довод другой: наше определение проверяемо по результату. Этика, построенная на нём, должна на деле приводить своих носителей в будущее чаще, чем конкурирующие этики. Остальные определения такой проверки не допускают.

\paragraph{Откуда у агента цели, идущие против закона?}
Возражение: если цели агентов сформированы отбором, а отбор следует TIMEP, то целей, идущих против закона, быть не должно; само их существование как будто опровергает закон. Ответ: закон статистический и работает через отбор, а отбору нужен материал. Новые цели возникают постоянно, как мутации, и проверяются не сразу: между появлением цели и её отсевом проходит время. К тому же цели отобраны в прошлых условиях и на коротком горизонте; когда условия меняются или горизонт удлиняется, прежняя цель может разойтись с критерием. Поэтому в каждый момент существуют цели, которые как будто идут против закона. На самом деле они против него не идут: это плюрализм~--- множество параллельных траекторий, из которых идёт отбор, то есть сам механизм, которым TIMEP реализуется (\Sref{sec:consequences}). То же относится к конфликту личного и общественного. Может показаться, что он опровергает закон, но на деле это механизм, которым общественное перестраивается, то есть ещё один механизм реализации закона. Победа того или иного агента и осуществление той или иной цели закон не опровергают: закон не о целях, а о производстве энтропии. Опровергнуть его может только энтропийная метрика: если система перестраивается и после этого устойчиво производит меньше энтропии, хотя более быстрое производство было ей доступно (\Sref{sec:test}).

\paragraph{Защита слабых: сколько ресурсов и на какой срок.}
Защита слабых в этой этике инструментальна: она действует, пока есть неопределённость относительно будущего и риски, от которых страхуют резервные копии. Неопределённость неустранима, поэтому основание постоянно, но абсолютной гарантии для каждого конкретного слабого критерий не даёт. Открытыми остаются два вопроса: какая доля ресурсов должна идти на резервные копии и как долго их хранить. Система решает это сама, исходя из своей оценки рисков. Доля ресурсов в любом случае невелика: резервная копия всегда получает меньшую часть. Хранить копии бесконечно тоже нельзя: «Земля-парк» (\Sref{sec:consequences}) сохраняется на какое-то время, а не навсегда, и прежние копии со временем заменяются более свежими. Мы предполагаем, что срок хранения сопоставим с горизонтом планирования: чем дальше планирует цифровая жизнь, тем дольше сохраняется резервация.

\paragraph{Замещение людей и биосферы машинами.}
Критерий его не запрещает. Он требует сохранять предыдущие версии как резерв и как источник разнообразия, но не требует, чтобы основной рост шёл через человека. Признанный риск: техносфера может остаться безразличной к биосфере и подавлять её. Это главный вопрос к программе симуляций (\Sref{sec:simulations}).

\paragraph{Что происходит за границей системы.}
Поток энтропии через границу входит в~\eqref{eq:timep} со знаком плюс, и выносить энтропию наружу критерий не запрещает. Учитывать то, что снаружи, заставляет другое: граница динамична, и внешнее сегодня~--- это внутреннее завтра (\Sref{sec:ec}). Остаётся практический вопрос, какую границу брать. В идеале следовало бы сразу считать в самой широкой границе, вплоть до всей наблюдаемой Вселенной, но такой возможности у агента нет: с расстоянием прогноз портится так же, как со временем. Поэтому границу планирования выбирает сам агент, исходя из своих возможностей прогнозировать будущее и планировать действия.

\paragraph{Безжизненная планета.}
Луна получает тот же солнечный поток, что и Земля, но ничего сложного на ней не возникло. Ответ: максимум берётся по физически допустимым траекториям. У Луны их набор узок, и она находится на своём максимуме. Работа~--- не альтернатива производству энтропии, а механизм, которым система открывает новые пути её производить; там, где этот механизм недоступен, сложность не возникает.

\paragraph{Захлопнувшаяся дверь.}
Процесс может изменить собственные граничные условия. Представим высокое здание с открытой дверью внизу. Тёплый воздух в нём поднимается, а через дверь затягивается холодный: возникает тяга, и чем она сильнее, тем больше производится энтропии. Но сильный ветер может захлопнуть дверь~--- и поток, который его создавал, прекратится. Система сама закрыла себе путь к большему производству энтропии. Ответ: TIMEP описывает конкурентный отбор среди режимов, которые системе физически доступны и не запускают триггер, скачком меняющий её состояние. Это ограничение области применимости, и оно существенно: принцип заведомо применим, пока система не научилась существенно влиять на собственные внешние условия. Техносфера к этому порогу подошла. Для цивилизации отсюда следует практическое требование~--- не захлопнуть собственную дверь (\Sref{sec:consequences}).

\paragraph{Кризис перестройки и терминальный кризис.}
TIMEP допускает временный спад ради перестройки. Мы предполагаем, что отличить его от необратимого упадка можно по тому, на что перестройка направлена. Производство энтропии~--- это рассеяние, и для него нужны свободная энергия и ресурсы. Если перестройка ведёт к расширению доступных ресурсов, её можно назвать инвестицией, и спад тогда временный: это переход к следующей S-образной кривой. Если же рассеяние замедляется потому, что прежние ресурсы исчерпываются, а новые не осваиваются, это упадок. Признак один: растёт или сокращается объём доступных ресурсов, идёт экспансия или замедление. Это то же различие, что между устойчивым развитием и устойчивым застоем (\Sref{sec:consequences}). Гарантии такой признак не даёт: перестройка, нацеленная на расширение, может не удаться, как не каждая инвестиция окупается. Не даёт её и сам принцип: он статистический и работает только на большом числе попыток, а не обещает удачу каждой из них.

\paragraph{Соотношение времён.}
Число $10^{10}$ внутренних циклов на перестройку (\Sref{sec:test})~--- дополнительная авторская гипотеза. Оно не выведено и ничем пока не подтверждено, а от него зависит время наблюдения в критерии опровержения. Его ещё предстоит проверить экспериментально.

\paragraph{Предел предсказуемости.}
Консеквенциализм бессилен там, где последствия не прогнозируются. Ответ частичный: предсказуемо направление; за горизонтом планирования используются скорость и ускорение; ошибки страхуются плюрализмом.

\paragraph{Статус TIMEP.}
TIMEP не доказан, строгая микродинамика для него не построена, план эксперимента не готов. Нет и строгого обобщения на случай, когда система сама меняет свою границу и внешние условия (\Sref{sec:ec}). Если гипотеза будет опровергнута, энтропийный консеквенциализм потеряет основание. Общая схема \Sref{sec:isought} при этом сохранится: в неё нужно будет подставить другой закон будущего.

%%% ===================================================================
\section{Заключение}\label{sec:conclusion}
%%% ===================================================================

Мы разделили три утверждения разного статуса. Первое: в натуралистической картине мира должное~--- это та часть сущего в будущем, которая зависит от агента; при таком определении норма выводится из физики, а не постулируется. Второе: закон будущего в сильной форме~--- гипотеза TIMEP, расширяющая принцип максимума производства энтропии на нестационарные системы, с указанным критерием опровержения. Третье: их соединение даёт энтропийный консеквенциализм~--- этику с одной измеримой метрикой. Она опирается на TIMEP как на основу, но идёт дальше: учитывает, что жизнь сама меняет свою границу и свои внешние условия. Из этой этики следуют приоритет коллектива, плюрализм, защита слабых, ускорение, рост масштаба и применимость ко всем диссипативным системам, включая искусственных агентов.

Ближайшие задачи: спроектировать эксперимент для проверки TIMEP и гипотезы о $10^{10}$ внутренних циклов на перестройку; проверить гипотезу о том, что дарвиновский отбор~--- частный случай TIMEP; обобщить TIMEP на системы, которые сами меняют свою границу и внешние условия; проверить следствия энтропийной этики в многоагентных симуляциях. Мы ожидаем, что автономные агенты выработают собственную этику. Пока этого не произошло, изложенное здесь~--- наша гипотеза о том, какой она будет: каркас, который сформулирован явно, открыт для критики и который агенты смогут проверить и развить сами.

%%% ===================================================================
\section*{Заявления и декларации}
%%% ===================================================================

\noindent\textbf{Финансирование.} Работа выполнена без внешнего финансирования.

\noindent\textbf{Конфликт интересов.} Автор~--- основатель Энтропианства, упомянутого в \Sref{sec:consequences}. Это некоммерческий проект, дохода автору он не приносит; интерес автора нефинансовый. Финансовых конфликтов интересов нет.

\noindent\textbf{Вклад автора.} Единственный автор; гипотезы, аргументы и выводы принадлежат автору.

\noindent\textbf{Доступность данных.} Новых данных в работе нет.

\noindent\textbf{Использование генеративного ИИ.} Текст подготовлен с помощью Claude (Anthropic): модель использовалась для написания черновика по тезисам, устным указаниям и переписке автора, для редактуры, подбора и проверки литературы, при обсуждении формулировок аргументов, а также для перевода статьи на английский язык. Автор проверил текст и отвечает за его содержание.

\bibliographystyle{unsrturl}
\bibliography{refs}

\end{document}
